\documentclass{article}

\usepackage{fancyhdr}
\usepackage{enumitem}
\usepackage{extramarks}
\usepackage{amsmath}
\usepackage{amsthm}
\usepackage{amsfonts}
\usepackage{tikz}
\usepackage[plain]{algorithm}
\usepackage{algpseudocode}
\usepackage{fontawesome}
\usetikzlibrary{automata,positioning}

%
% Basic Document Settings
%

\topmargin=-0.45in
\evensidemargin=0in
\oddsidemargin=0in
\textwidth=6.5in
\textheight=9.0in
\headsep=0.25in

\linespread{1.1}

\pagestyle{fancy}
\lhead{\hmwkAuthorName}
\chead{\hmwkClass\ (\hmwkClassInstructor\ \hmwkClassTime): \hmwkTitle}
\rhead{\firstxmark}
\lfoot{\lastxmark}
\cfoot{\thepage}

\renewcommand\headrulewidth{0.4pt}
\renewcommand\footrulewidth{0.4pt}

\setlength\parindent{0pt}

%
% Create Problem Sections
%

\newcommand{\enterProblemHeader}[1]{
    \nobreak\extramarks{}{Problem \arabic{#1} continued on next page\ldots}\nobreak{}
    \nobreak\extramarks{Problem \arabic{#1} (continued)}{Problem \arabic{#1} continued on next page\ldots}\nobreak{}
}

\newcommand{\exitProblemHeader}[1]{
    \nobreak\extramarks{Problem \arabic{#1} (continued)}{Problem \arabic{#1} continued on next page\ldots}\nobreak{}
    \stepcounter{#1}
    \nobreak\extramarks{Problem \arabic{#1}}{}\nobreak{}
}

\setcounter{secnumdepth}{0}
\newcounter{partCounter}
\newcounter{homeworkProblemCounter}
\setcounter{homeworkProblemCounter}{1}
\nobreak\extramarks{Spring 2026 %\arabic{homeworkProblemCounter}
}{}\nobreak{}

%
% Homework Problem Environment
%
% This environment takes an optional argument. When given, it will adjust the
% problem counter. This is useful for when the problems given for your
% assignment aren't sequential. See the last 3 problems of this template for an
% example.
%
\newenvironment{homeworkProblem}[1][-1]{
    \ifnum#1>0
        \setcounter{homeworkProblemCounter}{#1}
    \fi
    \section{Problem \arabic{homeworkProblemCounter}}
    \setcounter{partCounter}{1}
    \enterProblemHeader{homeworkProblemCounter}
}{
    \exitProblemHeader{homeworkProblemCounter}
}

%
% Homework Details
%   - Title
%   - Due date
%   - Class
%   - Section/Time
%   - Instructor
%   - Author
%

\newcommand{\hmwkTitle}{Homework 1}
\newcommand{\hmwkDueDate}{Thursday, January 29, 11:59pm via Gradescope.} %Late submission with $10\%$ penalty until Wednesday, January 24, 11:59pm.}
\newcommand{\hmwkClass}{CS 6515}
\newcommand{\hmwkClassTime}{}
\newcommand{\hmwkClassInstructor}{G. Brito}
\newcommand{\hmwkAuthorName}{Your Name Here}

%
% Title Page
%

\title{
    \vspace{2in}
    \textmd{\textbf{\hmwkClass:\ \hmwkTitle}}\\
    \normalsize\vspace{0.1in}\small{Due\ on\ \hmwkDueDate}\\
    \vspace{0.1in}\large{\textit{\hmwkClassInstructor\ \hmwkClassTime}}
    \vspace{3in}
}

\author{\hmwkAuthorName}
\date{}

\renewcommand{\part}[1]{\textbf{\large Part \Alph{partCounter}}\stepcounter{partCounter}\\}

%
% Various Helper Commands
%

% Useful for algorithms
\newcommand{\alg}[1]{\textsc{\bfseries \footnotesize #1}}

% For derivatives
\newcommand{\deriv}[1]{\frac{\mathrm{d}}{\mathrm{d}x} (#1)}

% For partial derivatives
\newcommand{\pderiv}[2]{\frac{\partial}{\partial #1} (#2)}

% Integral dx
\newcommand{\dx}{\mathrm{d}x}

% Alias for the Solution section header
\newcommand{\solution}{\textbf{\large Solution}}

% Probability commands: Expectation, Variance, Covariance, Bias
\newcommand{\E}{\mathrm{E}}
\newcommand{\Var}{\mathrm{Var}}
\newcommand{\Cov}{\mathrm{Cov}}
\newcommand{\Bias}{\mathrm{Bias}}

\begin{document}

\maketitle

\pagebreak

\begin{center}
\underline{\textbf{Instructions.}}
\end{center}

\noindent \texttt{For the graded problems, when asked to design an algorithm you must describe it in plain English. Do not use pseudocode! Your solution must include:
\begin{itemize}
    \item A description of your {\it divide step:} how many subproblems you have, what are the input of those? What is the criteria to recurse? 
    \item A description of your {\it conquering step:} how do you get your solution from the subproblems? What extra work do you need to perform?
    \item The {\it base case:} what rule stops your algorithm from recursing forever?
    \item The runtime (in big-O notation) of your design. You must justify your answer, and you may use The Master Theorem if you wish. Fastest (and correct) algorithms are worth more credit.
\end{itemize}
}

\begin{center}
\underline{\textbf{Suggested reading.}}
\end{center}

\texttt{From Algorithms by Dasgupta, Papadumitriou, and Vazirani, ([DPV] from now on) chpater 2.2 covers Recurrence Relations and The Master theorem, chapter 2.6 covers FFT. Chapter 2.3 covers MergeSort, which will be useful in this assignment.}

\vspace{1cm}


\begin{center}
 \faHandODown \faHandODown \faHandODown \\
Type your solutions in $\LaTeX$ or your favorite editor. Handwritten responses won't be accepted.\\
   \faHandOUp \faHandOUp \faHandOUp
\end{center}


\pagebreak


\begin{homeworkProblem}
A school maintains two lists of student IDs:

\begin{itemize}
    \item \textbf{List A:} IDs of all students currently enrolled in extracurricular clubs (size $n$).
    \item \textbf{List B:} IDs of students who have graduated (size $m$).
\end{itemize}

\noindent
The school wants to update the club lists by removing any students who have graduated. Design an efficient algorithm to produce a new list of students who are still active in clubs.\\

{\it Hint: you can use {\tt MergeSort} as a black-box.}

\subsection{Answer}
\textbf{Divide Step:}
The algorithm creates two subproblems by recursively sorting each input list:

\begin{itemize}
    \item First subproblem: Sort list A (enrolled students) using MergeSort 
    \item Second subproblem: Sort list B (graduated students) using MergeSort
\end{itemize}

These are independent sorting operations whose input are unsorted List A and List B respectively. The criteria to recurse is inherent to MergeSort i.e we recurse as long as the list being sorted has more than one element.

\textbf{Conquer Step:}
Once both lists are sorted, we merge the results by performing a two-pointer linear scan:

\begin{itemize}
    \item Initialize two pointers: one at the start of sorted list A, one at the start of sorted list B
    \item Initialize an empty result list
    \item Compare the student IDs at both pointers:
        \begin{itemize}
            \item If A's current ID equals B's current ID, this student has graduated, so skip them (advance both pointers)
            \item If A's current ID is less than B's current ID, this student is still active, so add them to the result list and advance A's pointer
            \item If A's current ID is greater than B's current ID, advance B's pointer (this graduated student wasn't in clubs)
        \end{itemize}
    \item Continue until we've processed all of list A
    \item Any remaining students in A who haven't been checked are added to the result list
\end{itemize}

The extra work needed is this single linear pass through both sorted lists, which takes O(n + m) time. \newline


\textbf{Base Case:}
The base case comes from MergeSort: when a list has only one element (or is empty), it is already sorted and we return it immediately without further recursion. \newline

\textbf{Runtime Analysis:}
Let's analyze each component:

\begin{itemize}
    \item Sorting list A: O(n log n) using MergeSort
    \item Sorting list B: O(m log m) using MergeSort
    \item Merging step: O(n + m) for the two-pointer scan
\end{itemize}

\textbf{Total runtime}: O(n log n + m log m)

\textbf{Justification}: The dominant terms are the two sorting operations. The linear merge is absorbed by the sorting costs since n log n grows faster than n for n > 1. We cannot use the Master Theorem here because this isn't a single divide-and-conquer recurrence—we're combining two separate sorting operations followed by a linear merge. The overall complexity is the sum of sorting both lists.

If n and m are comparable in size, we can express this as O((n + m) log(n + m)) as an upper bound, but the more precise analysis is O(n log n + m log m).

\end{homeworkProblem}

\begin{homeworkProblem}
You are given an array $A$ of $n$ distinct integers. You are also given the fact that there are exactly $k$ local maxima in this array. Devise an algorithm to find the global maxima of the array in $O(k*\log n)$ time. \textcolor{red}{Removed.}
\end{homeworkProblem}

\begin{homeworkProblem}
    You are given $n$ polynomials $p_1, p_2,\ldots, p_n$, each of degree at most $d$. Design an efficient algorithm to compute the product $p_1p_2\ldots p_n$.

\subsection{Answer}
\textbf{Divide Step:}
We divide the n polynomials into two subproblems:

\begin{itemize}
    \item Left subproblem: Compute the product of the first half of polynomials: $p_1, p_2,\ldots, p_{n/2}$
    \item Right subproblem: Compute the product of the second half of polynomials: $p_{n/2+1}, \ldots, p_n$
\end{itemize}

The input to each subproblem is half the number of polynomials. The criteria to recurse is that we have more than one polynomial to multiply. \\

\textbf{Conquer Step:}
After recursively computing:

\begin{itemize}
    \item Left result: L(x) = $p_1, p_2,\ldots, p_{n/2}$ (a polynomial of degree at most d·n/2)
    \item Right result: R(x) = $p_{n/2+1}, \ldots, p_n$ (a polynomial of degree at most d·n/2)
\end{itemize}

We need to compute their product: L(x) · R(x)
The extra work needed is multiplying two polynomials using the Fast Fourier Transform (FFT):

\begin{itemize}
    \item $L(x)$ has degree at most $d \cdot n/2$
    \item $R(x)$ has degree at most $d \cdot n/2$
    \item Their product has degree at most $d \cdot n$
    \item Using FFT, we can multiply these two polynomials in $O(dn \log(dn))$ time
\end{itemize}

\textbf{Base Case:}
When we have only one polynomial (n = 1):

\begin{itemize}
    \item Simply return that polynomial $p_1$
    \item No computation needed: O(1) time
\end{itemize}

\textbf{Runtime Analysis:}
Let T(n) be the time to multiply n polynomials, each of degree at most d. 

\textbf{Recurrence relation:} $T(n) = 2T(n/2) + O(dn \log(dn))$

At each level of recursion:

\begin{itemize}
    \item We split into 2 subproblems of size $n/2$
    \item The conquer step multiplies two polynomials of degree $\approx d \cdot n/2$ each, which costs $O(dn\log(dn))$ using FFT
\end{itemize}
\textbf{Using the Master Theorem:}

We have: $T(n) = 2T(n/2) + f(n)$ where $f(n) = O(dn \log(dn))$

Compare $n^{\log_2 2} = n^1 = n$ with $f(n) = dn \log(dn)$

Since $dn \log(dn)$ is larger than just $n$ (grows as $n \log n$), we're in Case 3 of the Master Theorem.

However, we need to be more careful here...

\textbf{More careful analysis:}
Actually, the recurrence is more precisely:

At the root level: we multiply two polynomials of degree $\approx dn/2$ each $\rightarrow O(dn \log(dn))$

At the next level: we have 2 problems, each multiplying polynomials of degree $\approx dn/4 \rightarrow 2 \cdot O(dn/2 \log(dn/2)) = O(dn \log(dn))$

At each of $\log n$ levels: total work is $O(dn \log(dn))$

\textbf{Total runtime:} $O(dn \log(dn) \cdot \log n) = O(dn \log n \log(dn))$

Since $d$ is constant with respect to $n$, we can also write this as $O(dn \log^2 n)$ when $d$ is fixed.

    
\end{homeworkProblem}

\begin{homeworkProblem}
    Georgia Tech is upgrading Wi-Fi along Techwood Drive, modeled as a straight line. There are $n$ dorms at integer mile-markers $P = [p_1, p_2, \dots, p_n]$ (not necessarily sorted). You may install at most $k$ identical routers. Each router can be placed \emph{anywhere on the line} (not necessarily at a dorm) and covers all dorms within distance $r$ of its location. That is, a router at position $x$ covers every dorm $p_i$ with $|p_i - x| \le r$.
\\

\noindent Design an efficient divide-and-conquer algorithm to \emph{minimize the necessary coverage radius $r$} so that all dorms are covered using at most $k$ routers.
\end{homeworkProblem}
\end{document}
